\documentclass[10pt,a4paper]{amsart}
\usepackage[margin=19mm]{geometry}
\usepackage{amsmath,amssymb,mathtools}
\usepackage[scr=rsfs,cal=euler]{mathalfa}
\usepackage{xfrac}
\usepackage[unicode,hidelinks,bookmarks=false]{hyperref}
\newcommand{\ie}{i.e.\ }
\newcommand{\cf}{cf.\ }
\newcommand{\resp}{respectively\ }
\newcommand{\Subsection}{\subsection}
\providecommand{\nobreakdash}{\nobreak\hbox{-}}
\setlength{\emergencystretch}{2em}
\multlinegap0pt
\usepackage{amssymb}
\usepackage{xcolor}
\newtheorem{proposition}{Proposition}
\newcommand{\vect}[1]{\boldsymbol{\mathbf{#1}}}
\title{\LARGE \bf Convex Framework for Constrained Quadcopter Trajectory Optimization and Tracking with Continuous-Time Guarantees}
\author{Victor Freire, Xiangru Xu}
\date{Source: arXiv:2111.00951v1 (CC BY 4.0)}
\begin{document}
\maketitle

\noindent\emph{Source and licence.} \LARGE \bf Convex Framework for Constrained Quadcopter Trajectory Optimization and Tracking with Continuous-Time Guarantees. Version arXiv:2111.00951v1 (CC BY 4.0), distributed by arXiv under the Creative Commons Attribution 4.0 International licence, http://creativecommons.org/licenses/by/4.0/, as recorded in the arXiv licence metadata for this exact version and shown on the abstract page https://arxiv.org/abs/2111.00951v1. Full licence notice: \texttt{licenses/arxiv-2111.00951-CC-BY-4.0.txt}.\par\smallskip

\noindent\textit{Editorial scope.} Counted unit: the article's proposition convex\_sets\_prop (source0.tex lines 1063--1069). The article's complete original proof of it is delivered with it (source0.tex lines 1070--1071, one \texttt{proof} environment of 2 lines). No mathematics is written, repaired or re-derived: the statement and the proof are the article's own lines, in the article's own notation, and no retained line is substituted or re-flowed.  Retained uncounted as the source-local context the proof needs: lines 253--256; lines 259--262; lines 505--519.  Nothing was omitted from any retained span, so the package carries no omission record.  No retained line contains a citation, so the wrapper carries no bibliography and no bibliography entry.  No retained line contains a figure environment, an image-inclusion command or drawing code, so no image is delivered and the package declares no asset dependency.  Editorial formatting only, in addition: an AMS \texttt{amsart} preamble that restates the article's own declarations and macros used by the retained lines, namely the environments \texttt{proposition} on separate counters, together with the standard packages those lines need.  The retained environments print in this document as Proposition 1 in order of appearance, whereas the article prints the same items with the numbering of its own full document.  The complete native source of the article as distributed in the arXiv e-print for this exact version is bundled unchanged as \texttt{sources/117937/source0.tex}.
\begin{align}
    \text{(clamped)} \quad &\tau_0=\ldots=\tau_d, \quad \tau_{v-d}=\ldots=\tau_{v},\label{clamped}\\
    \text{(uniform)}\quad &\tau_{d+1}-\tau_{d}=\ldots =\tau_{N+1}-\tau_{N},\label{uniform}
\end{align}

\begin{equation}
    \vect{s}^{(r)}(t) = \sum_{i=0}^{N}\vect{P}_i\vect{b}_{r,i+1}^T\vect{\Lambda}_{d-r}(t) = PB_r\vect{\Lambda}_{d-r}(t),
    \label{b-spline_curve}
\end{equation}

\begin{proposition} \label{inclusion_prop}
Given a convex set $\mathcal{S}$ and the $r$-th derivative of a clamped B-spline curve $\vect{s}^{(r)}(t)$ defined as in \eqref{b-spline_curve}, if
\begin{equation}
     %\mbox{i)}\quad \quad
     \vect{P}_{j}^{(r)} \in \mathcal{S}, \quad j=r,\ldots,N
    \label{inclusion_global}
\end{equation}
holds, then $\vect{s}^{(r)}(t) \in \mathcal{S}, \ t\in[\tau_0,\tau_v)$. If
\begin{equation}
    %\mbox{ii)}\quad\quad
    \vect{P}^{(r)}_j \in S, \quad j=i-d+r,\dots,i, \quad i\in\{d,\ldots,N\}
    \label{inclusion_local}
\end{equation}
holds, then $\vect{s}^{(r)}(t)\in \mathcal{S}, \ t\in[\tau_i,\tau_{i+1})$.
\end{proposition}

\begin{proposition} \label{convex_sets_prop}
Suppose that there exist $n_s$ overlapping convex sets $\{\mathcal{S}_1,\ldots,\mathcal{S}_{n_s}\}$ inside the safe flight space, i.e., $\mathcal{S}_l\subset \mathcal{F},\ l=1,\ldots,n_s,$ satisfying $\mathcal{S}_i\cap \mathcal{S}_{i+1}\neq \emptyset, \ i = 1,\ldots,n_s-1$ and consider B-spline curve $\vect{r}(t)$ defined as in \eqref{b-spline_curve} with $N=n_s+d-1$. If the following conditions hold:
\begin{equation}\label{eqobstc}
    \vect{P}_{j-1} \in \mathcal{S}_{l},\quad j = l,\ldots,l+d, \quad l = 1,\ldots,n_s,
\end{equation}
then  $\vect{r}(t)\in\mathcal{F}$  for all $t\in[\tau_0,\tau_v)$ and $\vect{r}(t)$ is at least $C^{d-1},\  \forall t\in[\tau_0,\tau_v)$.
\end{proposition}

\begin{proof} Applying Proposition \ref{inclusion_prop} to the inclusion constraint of the $l$-th  group of control points results in $\vect{r}(t)\in \mathcal{S}_l, \ t\in[\tau_{l+d-1},\tau_{l+d})$, which means that the $l$-th segment of the B-spline curve $\vect{r}(t)$ is contained within $\mathcal{S}_l$. Taking now $l=1,\ldots,n_s$ and recalling that we consider a clamped knot vector \eqref{clamped} results in $\vect{r}(t)\in\mathcal{F},\ t\in[\tau_0,\tau_v)$. Further, note that $\mathcal{S}_l\cap \mathcal{S}_{l+1} \neq \emptyset$ allows $\vect{r}(t)$ to be continuous despite the segment-wise constraints. 
\end{proof}

\end{document}
